\section{A height bound from odd contact with five points}\label{hei:section}

We prove Theorem~\ref{hei:height}, whose contact radical supplies the
primes used in Section~\ref{sec:reduction}.  For
\(s\in\mathbf P^1(\Q)\setminus\mathcal C\), the target is
\[
                         \hgt(s)\le H M_S(s)^c,
\]
where \(M_S(s)\) is the squarefree product of the nonexceptional primes
of odd contact with the five fixed points, as defined in that theorem.

Throughout this section, $\hgt$ is the absolute logarithmic Weil height,
and $h_{\mathrm F}$ is the original stable Faltings height.  Unless
another dependence is stated, constants depend only on the geometric
objects and auxiliary levels fixed in the proof.

The geometric part of the proof constructs a finite cover
$\pi:Y\to\mathbf P^1$ with ramification index two exactly above
$\mathcal C$, carrying a family of principally polarized abelian surfaces.
Outside fixed exceptional primes, the local equation becomes $t=z^2$
after a finite unramified extension.  Thus only odd contact can ramify
the splitting field of a rational fiber, as Lemma~\ref{hei:fiber-field}
makes precise.
The polarized isomorphism class of the surfaces varies on each component,
which also allows their Faltings height to control the base height
(Lemma~\ref{hei:base-height}).

The surfaces carry quaternionic multiplication, and fibers above a common
base point are isogenous.  For a rational base point, this gives isogenies
between a surface and its Galois conjugates.  Fibers with extra
endomorphisms have uniformly bounded height by complex multiplication.
For the other fibers, a controlled splitting of the isogeny obstruction
produces a two-dimensional Galois representation.  Modularity and a
conductor estimate give a weight-two newform whose level is bounded by a
fixed power of $M_S(s)$.  The surface is then an isogeny factor of the
corresponding modular Jacobian over a field of degree $O(M_S(s))$.
Height bounds for that Jacobian and the quantitative isogeny theorem
complete the estimate.


\subsection{The five-point quaternionic quotient}

We first construct the cover whose ramification will detect the odd
contact orders in Theorem~\ref{hei:height}.  Let $B/\Q$ be the
indefinite quaternion algebra of discriminant
$210=2\cdot3\cdot5\cdot7$, let $\mathcal O$ be a maximal order, and put
$G=B^\times$.  Fix $B\otimes_\Q\mathbf R\simeq M_2(\mathbf R)$.
The quaternionic Shimura datum is $(G,\mathfrak H^\pm)$, where the
two half planes parametrize the conjugates of the elliptic-curve
Hodge homomorphism.  Its reflex field is $\Q$: over $\mathbf C$ its
cocharacter has the usual $\mathrm{GL}_2$ conjugacy class, and this
class is invariant under the Galois action for the inner form $G$.

We use canonical models and functoriality for Shimura varieties
\cite{Deligne1971,Deligne1979}.  In the form needed here, they give
the model over the reflex field with complex points
\[
 G(\Q)\backslash
       \bigl(\mathfrak H^\pm\times G(\mathbf A_f)/K\bigr)
\]
for a compact open subgroup $K$, and define level maps and
finite-adelic right actions over that field.  We use the full Shimura
varieties over $\Q$, including all their geometric components.

At maximal level, put
\[
 X=\operatorname{Sh}_{\widehat{\mathcal O}^{\,\times}}
                         (G,\mathfrak H^\pm).
\]
This curve is proper because $B$ is a division algebra.  Strong
approximation for $B^1$, the reduced-norm theorem, and
$\operatorname{Nrd}(\mathcal O_p^\times)=\Z_p^\times$ give
\[
 G(\mathbf A_f)=G(\Q)^+\widehat{\mathcal O}^{\,\times}.
\]
Indeed, match the finite-adelic norm modulo $\widehat\Z^\times$ by
a positive rational reduced norm and then apply strong approximation
to the norm-one part.  There are also rational elements of negative
norm, since no real place is ramified.  Consequently $X_{\mathbf C}$
is the connected compact curve
\[
       (\mathcal O^1/\{\pm1\})\backslash\mathfrak H.
\]

The uniformizing group has no nontrivial elliptic stabilizers.  A
noncentral norm-one unit fixing a point has integral trace of absolute
value less than two, hence trace $0$ or $\pm1$.  It would generate
$\Q(i)$ or $\Q(\sqrt{-3})$.  These fields cannot embed in $B$, since
they split at the ramified primes $5$ and $7$, respectively.  The
Eichler area formula, in the normalization of
\cite[Theorem 39.1.2]{Voight2021}, therefore gives
\[
 \frac{\operatorname{area}(X(\mathbf C))}{2\pi}
       =\frac16\prod_{p\mid210}(p-1)=8.
\]
There are no cusps or elliptic corrections, so $2g(X)-2=8$ and
$g(X)=5$.

The finite-adelic normalizer of $\widehat{\mathcal O}$ induces the
Atkin--Lehner group
\[
                  \mathcal W\simeq(\Z/2\Z)^4
\]
over $\Q$.  At a split prime, the normalizer is scalars times order
units.  At a division prime, its quotient by scalars and order units
has order two, detected by parity of reduced-norm valuation.  Finite
adelic scalars act trivially, since
$\mathbf A_f^\times=\Q^\times\widehat\Z^\times$.  The resulting
action is faithful: a generic domain point has rational stabilizer
equal to the center, which accounts precisely for the scalars already
removed.  Let $X^*=X/\mathcal W$.

The geometry of this discriminant-\(210\) quotient has been tabulated
before.  Long, Maclachlan, and Reid list the signature
\((0;2^5;0)\), recording a genus-zero quotient with five elliptic
cycles of order two and no cusps
\cite[Table~3]{LongMaclachlanReid2006}.  Nualart Riera gives genus
\(5\) for the original curve and the five involution labels
\(30,42,70,105,210\), and identifies the full quotient over
\(\mathbb Q\) with \(\mathbf P^1\)
\cite[Propositions~4.1--4.2]{Nualart2015}.  The calculation here
recovers this fixed configuration and establishes the rational branch
values used in the reduction.

A finite stabilizer of a point on a smooth characteristic-zero curve
acts faithfully on its tangent line and is cyclic.  Since every
nonidentity element of $\mathcal W$ has order two, each nontrivial
stabilizer has order two.  If $r$ is the number of geometric branch
values of $X\to X^*$, Riemann--Hurwitz gives
\begin{equation}\label{hei:hurwitz}
             2g(X^*)-2+\frac r2=\frac{2g(X)-2}{16}=\frac12.
\end{equation}
Thus $r=5-4g(X^*)\le5$.

We exhibit five different stabilizer labels.  Take
\[
                         m=30,42,70,105,210.
\]
The field $\Q(\sqrt{-m})$ is nonsplit at every prime dividing $210$.
For the first three values, the only prime of $210$ not dividing $m$
is respectively $7,5,3$, and the residues of $-m$ there are the
nonsquares $5,3,2$.  For $m=105$, the quadratic discriminant is
$-420$, so $2$ is ramified too.  For $m=210$, all four primes are
ramified.  The quaternion embedding criterion from
Albert--Brauer--Hasse--Noether, as in \cite{Voight2021}, gives an
embedding of each field in $B$.

An embedded $\sqrt{-m}$ lies in some maximal order: at split local
places an integral element stabilizes a lattice, and at division
places it lies in the unique maximal order.  Intersecting these local
orders gives a global maximal order containing the embedded element.
Strong approximation implies that all maximal orders in this
indefinite rational quaternion algebra are conjugate by $B^\times$
\cite[Theorems 28.2.10 and 28.2.11(b)]{Voight2021}, so a conjugate
of the embedded element lies in our fixed $\mathcal O$.
It normalizes $\mathcal O$ because it is a unit at every prime outside
$210$, and at a division prime every element normalizes the unique
maximal order.  Its norm is $m$ and it fixes a point of $\mathfrak H$,
so it gives a fixed point of the involution $w_m$ with the corresponding
norm-parity label.

The five labels are different.  A fiber of the quotient is one
$\mathcal W$-orbit, and the stabilizers along this orbit agree because
$\mathcal W$ is abelian.  Distinct labels thus yield distinct branch
values, proving $r\ge5$.  Equation~\eqref{hei:hurwitz} now forces
$r=5$ and $g(X^*)=0$.  There is exactly one branch value for each of
the five labels.  Every $w_m$ is defined over $\Q$, so Galois preserves
its uniquely labelled branch value.  Each branch value is therefore
individually rational.  In particular $X^*$ is a genus-zero curve
with a rational point, hence is isomorphic over $\Q$ to
$\mathbf P^1$.  Fix an isomorphism sending one branch value to
$\infty$, and denote the five branch values by $\mathcal C$.

\newpage
\subsection{The abelian surface family}

We next pass to a fine level so that the cover carries abelian
surfaces.  The construction will provide a finite level map $Y\to X$,
unramified over $X$, and a principally polarized abelian surface
scheme $\mathcal A\to Y$.  Write $\pi$ for the composite
\[
 Y\longrightarrow X\longrightarrow X^*\simeq\mathbf P^1,
 \qquad \pi:Y\longrightarrow\mathbf P^1.
\]
Every fiber $\mathcal A_y$ with $y\in Y(\overline\Q)$ will have an embedding
$B\hookrightarrow\operatorname{End}^0_{\overline\Q}(\mathcal A_y)$
given by left multiplication.
Fibers with the same image under $\pi$ will be geometrically isogenous,
and the polarized moduli map will be nonconstant on every geometric
component.  We now construct the family and prove these properties.

Canonical-model functoriality associates algebraic maps to morphisms
of Shimura data at compatible levels
\cite{Deligne1971,Deligne1979}.  A Siegel variety at principal level
at least three carries the universal principally polarized abelian
scheme.  The required morphism to a Siegel datum comes from the
following symplectic representation.

On the four-dimensional rational space $V=B$, put
\[
 \rho(g)(x)=x\bar g,\qquad
 \psi(x,y)=\operatorname{Trd}(\Delta x\bar y),
\]
where the bar is the standard quaternion involution and $\Delta$ is a
pure quaternion with $\Delta^2<0$.  Such a rational $\Delta$ exists
by real approximation in the trace-zero subspace.  The identities
$\bar\Delta=-\Delta$ and cyclicity of reduced trace show that $\psi$
is alternating.  It is nondegenerate because the reduced-trace
pairing is nondegenerate and $\Delta$ is invertible.  Moreover
\[
 \psi(x\bar g,y\bar g)=\operatorname{Nrd}(g)\psi(x,y).
\]
Quaternion conjugation reverses products, as does composition of
right multiplications; together these identities give
$\rho(gh)=\rho(g)\rho(h)$.
Thus $\rho$ is a faithful symplectic-similitude representation.

The Hodge homomorphism makes $V_{\mathbf R}$ a complex vector space by
right multiplication.  At a base point represented by
$J=\left(\begin{smallmatrix}0&-1\\1&0\end{smallmatrix}\right)$,
the associated real quadratic form is
\[
 \psi(x,x\bar J)=\operatorname{tr}(J\Delta xx^{\mathrm t}).
\]
Since $\Delta$ has trace zero, $J\Delta$ is symmetric, and its
determinant is positive; it is therefore definite.  Choose its sign
on one half plane.  Conjugation by $G(\mathbf R)$ gives all the other
complex structures, and the displayed similitude identity gives the
corresponding definite sign on the other half plane as well.  This is
the two-component Siegel datum.  A real scalar acts on $V$ with the
homological weight of an elliptic curve, which is the required weight.

Choose a rational symplectic basis for $V$.  The standard lattice in
that basis is self-dual for $\psi$.  Intersect a sufficiently small
compact open subgroup of $\widehat{\mathcal O}^{\,\times}$ with the
inverse image of a principal Siegel level, and call the resulting
subgroup $K_Y$.  Put
\[
                 Y=\operatorname{Sh}_{K_Y}(G,\mathfrak H^\pm).
\]
The induced morphism from $Y$ to the fine Siegel moduli scheme pulls
back its universal family to the principally polarized abelian
surface scheme $\mathcal A\to Y$.
The self-dual lattice need not be preserved by the whole maximal
order: the intersection of levels provides the integral moduli
problem, while the rational Hodge structure retains the $B$-action.
Fix this level and family once and for all.

There may be finitely many geometric components of the smooth
projective curve $Y$.  The level map $Y\to X$ is unramified by the
uniformization and the absence of elliptic stabilizers on $X$.
Consequently its composite $\pi$ has ramification index two at
every point above $\mathcal C$ and index one elsewhere.

Every geometric fiber $A=\mathcal A_y$ has left $B$-multiplication in
$\operatorname{End}^0_{\overline\Q}(A)$.  Left multiplication commutes
with the right complex structures on $V$; the equivalence between
abelian varieties up to isogeny and their polarizable rational Hodge
structures supplies the endomorphisms.  Homomorphisms of abelian
varieties defined over $\overline\Q$ do not acquire new elements after
extension to $\mathbf C$ \cite{Mumford1970}.

For non-CM points, the moduli construction of Guitart and Molina gives
a related description by compatible isogenies
\cite[Lemma~1, Corollary~2, and Section~3]{GuitartMolina2009}.
For this fine-level family, the following Hodge argument gives the
underlying geometric isogenies for all fibers needed here; compatible
choices are made in the descent argument.  Changing level or applying the normalizer changes
the finite-adelic coordinate, while domain coordinates of fibers
above a common point of $X^*$ can be identified up to $G(\Q)$.
Their rational Hodge structures are therefore isomorphic, so the
fibers are geometrically isogenous.  In particular, if
$\pi(y)=s\in\mathbf P^1(\Q)$, then $\pi({}^g y)=g(s)=s$ and
$\mathcal A_y$ is isogenous to each of its Galois conjugates.

Finally, the polarized moduli map of the family is nonconstant on
every geometric component.  Its complex structures vary in a half
plane, whereas a fixed rational Hodge structure has only countably
many presentations on the fixed rational space $V$.  Thus a component
cannot have a single geometric isomorphism class of fibers.

\subsection{Rational fibers and odd-contact ramification}

The ramification index two now explains the parity in the theorem.
We make precise how even contact with a branch value gives unramified
local lifts, and record the good reduction of the resulting surfaces.

Enlarge a fixed finite set $S$ of rational primes whenever necessary
in the following construction.  All these enlargements depend only
on the fixed curves, maps, and family, and occur before choosing $s$.
Include $2$, primes at which the branch sections meet, and primes at
which the chosen coordinates or the geometric models degenerate.
Spread $Y$, $\mathbf P^1$, $\pi$, and $\mathcal A$ over
$\Z[S^{-1}]$ so that the curves are smooth and proper, $\pi$ is finite
and étale away from the branch sections, and
\[
                         \pi^*b=2D_b
\]
as relative Cartier divisors, with $D_b$ the reduced inverse image,
finite étale over the ground ring.  The abelian scheme also spreads
over this proper model of $Y$.

\begin{lemma}\label{hei:fiber-field}
For $s\in\mathbf P^1(\Q)\setminus\mathcal C$, put $M=M_S(s)$.
There is a Galois extension $K_0/\Q$ of uniformly bounded degree
splitting the fiber $\pi^{-1}(s)$, unramified outside
$S\cup\{q:q\mid M\}$.  Every fiber $A=\mathcal A_y$ above $s$ is
defined over $K_0$ and has good reduction at every place of $K_0$
outside $S$.
\end{lemma}

\begin{proof}
The fiber has at most $\deg\pi$ geometric points, so its splitting
field is Galois of degree at most $(\deg\pi)!$.  Fix $q\notin S$.
If the reduction of $s$ avoids the branch sections, finite étaleness
implies that the fiber is unramified at $q$.  Otherwise it meets
exactly one branch section $b$, and its contact order is
\[
             v_q(L_b(u,v))-\min\{v_q(u),v_q(v)\}.
\]
After a finite unramified local extension, all points of $D_b$ over
this reduction are defined, since $D_b$ is finite étale.  At each
such point choose a formal coordinate $z$ cutting out $D_b$.
If $t$ cuts out $b$ on the base, the divisor equality gives
$\pi^*t=z^2e(z)$ with $e(z)$ a unit.  After a further unramified
extension its residue has a square root, and Hensel's lemma, since
$q\ne2$, gives a square root of $e(z)$ in the formal power-series
ring.  Replacing $z$ by $z\sqrt{e(z)}$ changes the local equation to
$t=z^2$.  Make this calculation at every point of $D_b$ over the
branch reduction, using a common further finite unramified extension.

If the contact order is even, a square root of $t(s)$ is obtained
over an unramified extension by extracting the square root of its
unit part.  Both local lifts at every such point are consequently
unramified.  By properness, every geometric point of the fiber specializes
to one of the points of $D_b$ just treated, so the whole fiber is
unramified.  This proves
that ramification can occur only at the primes in the assertion.

Every $K_0$-point $y$ extends over the local valuation rings outside
$S$ by properness of $Y$.  Pulling back the spread abelian scheme
along these sections gives good reduction of $\mathcal A_y$.
This also applies at places where $K_0/\Q$ is ramified.  Thus field
ramification may occur at primes dividing $M$, while good reduction
over $K_0$ holds outside the fixed set $S$.
\end{proof}

This use of contact multiplicities has a classical specialization
precedent.  Darmon and Granville work with a fixed Galois cover of
the projective line branched at \(0,1,\infty\).  In that setting,
their Proposition~3.2, credited there to Beckmann, relates
divisibility of contact orders by branch indices to the absence of
new ramification outside fixed bad primes in specialized fiber fields
\cite{DarmonGranville1995}.  For the cover used here, the local
equation \(t=z^2\) makes the parity of the contact order relevant.

\subsection{Comparison of base height and Faltings height}

The fiber-field lemma identifies the primes that may ramify when a rational
base point is lifted to the family.  We also need to recover the height
of that base point from the fiber.  The next lemma depends only on the
fixed family, and applies to every algebraic point of $Y$.

\begin{lemma}\label{hei:base-height}
For every $y\in Y(\overline\Q)$ one has
\begin{equation}\label{hei:comparison}
 \hgt(\pi(y))\le C\bigl(1+\max\{0,h_{\mathrm F}(\mathcal A_y)\}\bigr),
\end{equation}
with a constant independent of $y$ and its field of definition.
\end{lemma}

\begin{proof}
We will use a theta-null map with two properties: its projective height
agrees up to a fixed additive constant with Pazuki's theta height
$h_\Theta$ of the fiber, and the map is nonconstant on each component
of a fixed cover of $Y$.  Here $h_\Theta$ is the projective height of
the theta-null point for the corresponding polarization and level
\cite[Section 2.3]{Pazuki2012}.  The second property makes the pulled-back
$\mathcal O(1)$ ample on the covering curve.  Comparing that ample
height with the pullback of the base height gives
\[
 \hgt(\pi(y))\ll 1+h_\Theta(\mathcal A_y).
\]
Pazuki's comparison of theta and Faltings height will then prove the
lemma.  We first construct a map with the two stated properties.

For every fixed even integer $r\ge2$, the required algebraic theta
data can be chosen after a finite cover of each geometric component
of $Y$.  Over a finite extension of its function field, choose a
symmetric ample line bundle $\mathcal M$ representing the principal
polarization, together with the theta data of
\cite[Section 2.3]{Pazuki2012}.  These choices require only a finite
extension: a representative of a geometric principal polarization
can be made symmetric by translation, and the torsion points and
finite theta data are defined after a finite extension.  They spread
over a dense open subset of the smooth projective covering curve,
where the theta-null map is algebraic.  The value of this map at a
fiber is the zero-section point $\Theta(0)$ constructed from Pazuki's
good choices of theta data; after complex uniformization of that fiber,
\cite[Sections~2.3.2--2.3.4]{Pazuki2012} identifies this point, up to the finite
coordinate ambiguity among good choices, with the characteristic
Nullwerte displayed below.  These covers and data may
depend on $r$; we will fix the final level after proving nonconstancy.

We next compare the projective height of this map with Pazuki's theta
height, uniformly in the fiber at any fixed $r$.  For a fiber
$A=\mathcal A_y$ and a period matrix $\tau$ for its principal
polarization, the characteristic coordinates can be taken to be
\begin{equation}\label{hei:theta-series}
 \theta[a,b](\tau,0)=
 \sum_{n\in\Z^2}
   \exp\bigl(\pi i(n+a)^{\mathrm t}\tau(n+a)
                       +2\pi i(n+a)^{\mathrm t}b\bigr),
 \quad a,b\in(r^{-1}\Z^2)/\Z^2.
\end{equation}
Start with a nonzero section of $\mathcal M$, whose space of sections
has dimension one, pull it back by $[r]$, and use
$[r]^*\mathcal M\simeq\mathcal M^{r^2}$.  Translating by the theta
group over representatives of $A[r^2]/A[r]$ gives these coordinates.
Changing the generating section multiplies all coordinates by one
common scalar.  Two symmetric representatives differ by a
two-torsion element of $\operatorname{Pic}^0$, which shifts the
half-characteristic; these shifts are included because $r$ is even.

Put $k=r^2$.  At the generic fiber of the chosen cover, let $g_x$
denote the good theta-group lift of a $k$-torsion point $x$.  These
lifts satisfy the good-choice relation $g_xg_y=\zeta_{x,y}g_{x+y}$ with
$\zeta_{x,y}\in\mu_k$.  It gives $g_0\in\mu_k$ and, by iteration at
$kx=0$, $g_x^k\in\mu_k$; hence $g_x^{k^2}=1$.  Choose $u_x=c_xg_x$
with $u_x^{2k}=1$ by taking a $2k$-th root of the central scalar
$(g_x^{2k})^{-1}$.  Since $2k\mid k^2$, we have
$c_x^{k^2}=u_x^{k^2}g_x^{-k^2}=1$.  These finitely many roots can be
chosen once after a finite constant extension of the function field
and spread with the good data.  We use these normalizations only to
compare coordinates of the fixed good-choice theta-null map.

In analytic theta coordinates, a lift of translation by
$(m+\tau l)/k$ is given by shifting the argument and multiplying by
\[
          \exp(\pi i l^{\mathrm t}\tau l/k+2\pi i l^{\mathrm t}z).
\]
Theta automorphy shows that its order divides $2k$ and that it permutes
$\theta[a,b](\tau,rz)$, up to bounded-order roots of unity, by
$(a,b)\mapsto(a+l/r,b+m/r)$.  Two lifts of the same translation whose
$2k$-th powers are one differ by a scalar in $\mu_{2k}$.  Thus
comparison with the algebraic $u_x$ introduces a $\mu_{2k}$ factor
after the preceding $\mu_{k^2}$ change from the good lift.  All these
coordinate roots of unity preserve projective height, as do
permutations and a common scalar.

The customary Fourier change of theta basis is a fixed invertible
matrix at this level, so it changes height by $O_r(1)$.  Grouping
$r(n+a)$ modulo $k$ in~\eqref{hei:theta-series} relates these
coordinates blockwise by finite Fourier matrices to
$\theta[e/k,0](k\tau,0)$, $e\bmod k$.  There are $r^4$ independent
sections, as required by $h^0(\mathcal M^{r^2})=r^4$, and their
values are not all zero since $\mathcal M^{r^2}$ is basepoint-free.
Pazuki's exact theta height uses the $\ell^2$ projective height
\cite[Definition 2.6]{Pazuki2012}; its difference from the usual
projective height is also bounded at this fixed level.  These
comparisons give one additive constant depending on $r$, not on the
fiber or its field of definition.

It remains to choose a level at which the theta-null map is
nonconstant.  The following argument avoids an assumption about
generic injectivity of a particular theta map.  Take a connected
analytic disk with a fixed homology marking on which the period matrix
varies.  Suppose the projective vectors
in~\eqref{hei:theta-series} were constant at every sufficiently
divisible even level.  At a chosen period $\tau_0$, select rational
$b_0$ with $\theta[0,b_0](\tau_0,0)\ne0$; such a $b_0$ exists since
the function of real $b$ has constant Fourier coefficient one.
Shrink the disk so its denominator stays nonzero.  At a common level
containing $b$ and $b_0$, projective constancy implies that
\[
          \frac{\theta[0,b](\tau,0)}{\theta[0,b_0](\tau,0)}
\]
is constant as $\tau$ varies, for every rational $b$.  Density of
rational $b$, continuity, and absolute convergence of the theta
series show that the functions of real $b$ for $\tau$ and $\tau_0$
are proportional.  Their constant Fourier coefficients are one, so
they are equal.  Consequently every
$\exp(\pi i n^{\mathrm t}\tau n)$ is constant.  Taking $n=e_i$ and
$n=e_i+e_j$ and using continuity forces $\tau$ itself to be constant,
a contradiction.  A level with a nonconstant vector therefore exists.
Nonconstancy persists at multiples of that level because the old
characteristic coordinates occur among the new ones.  A common even
multiple works for the finitely many components.

Fix this common level, and make the preceding choices of covers and
theta data at that level.  The finitely many covering curves, maps,
and open subsets are now fixed, and can all be defined over a fixed
number field.  On each base component, choose a smaller marked disk
avoiding the branch values of the covers and the images of their
omitted points, and take a connected lift to each covering curve above
it.  The selected holomorphic characteristic vector remains
nonconstant on the smaller disk.  If the theta-null map were constant
on a covering curve, the preceding pointwise identification would
confine that vector on the connected lift to a fixed finite set.
Continuity and connectedness would then make it constant.  Thus the
theta-null map is nonconstant on each covering curve and extends
across the omitted points.  Its pullback of $\mathcal O(1)$
has positive degree and is ample.  A sufficiently large multiple
of this line bundle minus the pullback by $\pi$ of
$\mathcal O_{\mathbf P^1}(1)$ also has positive degree.  The height
machine and the lower bound for an ample height on a projective curve
therefore give
\begin{equation}\label{hei:theta-base}
                  \hgt(\pi(y))\le C(1+h_\Theta(\mathcal A_y))
\end{equation}
on the open where the theta data are defined
\cite{BombieriGubler2006}.

Here $h_\Theta$ is Pazuki's exact theta height; the bounded coordinate
and projective-height comparisons have already been absorbed
in~\eqref{hei:theta-base}.  Pazuki uses the Deligne normalization
\[
 h_{\mathrm F}^{\mathrm D}(A)
   =h_{\mathrm F}(A)+\frac{\dim A}{2}\log\pi;
\]
compare \cite[Section 2.2]{Pazuki2012} and
\cite[Section 2.3]{GaudronRemondPeriods2014}.  Here $h_{\mathrm F}$
is the original stable Faltings height.  For the present surfaces
the difference is the fixed constant $\log\pi$.

For $H_\Theta=\max\{h_\Theta,1\}$ and
$H_{\mathrm D}=\max\{h_{\mathrm F}^{\mathrm D},1\}$,
\cite[Corollary 1.3(2)]{Pazuki2012} gives
\[
 |H_\Theta-\tfrac12H_{\mathrm D}|
 \le C_2\log(\min\{H_\Theta,H_{\mathrm D}\}+2),
\]
where $C_2$ depends only on dimension and level.  Since
$\log(H_{\mathrm D}+2)\le2H_{\mathrm D}$ and
$H_{\mathrm D}\le\max\{h_{\mathrm F},1\}+\log\pi$ for these surfaces,
we obtain
\[
 h_\Theta\le C(1+\max\{0,h_{\mathrm F}\}).
\]
Combining this with~\eqref{hei:theta-base} proves the assertion on the
chosen open sets.

The omitted points on the fixed covering curves form a finite
geometric set.  Their base heights are bounded, and enlarging $C$
handles them.  Absolute heights and stable Faltings heights are
unchanged by field extension, so passing to points of the fixed
covers introduces no dependence on their residue-field degrees.
Taking the maximum over the finitely many components proves the lemma.
\end{proof}

\begin{remark}[An alternative deduction on the compact moduli image]
The same numerical comparison can alternatively be deduced by bounding
the period term in Pazuki's Theorem 1.1 for this family.  That theorem compares
$h_\Theta-\tfrac12h_{\mathrm F}^{\mathrm D}$ with the average of
logarithms of determinants of imaginary parts of Siegel-reduced
periods, with bounded error at fixed dimension and level
\cite{Pazuki2012}.  The determinant terms are bounded here.
The moduli image is compact because $Y$ is projective, and the
minimum nonzero period length in the polarization metric has a
positive lower bound on it.  On the subgroup $\Z^2$ of the period
lattice the metric has matrix $(\operatorname{Im}\tau)^{-1}$.
Minkowski's theorem bounds $\det\operatorname{Im}\tau$ above in terms
of that minimum length, and Siegel reduction bounds it below by a
positive constant.  These bounds hold at every embedding, because
every conjugate belongs to the fixed compact moduli image.  Theorem 1.1
and the same normalization conversion therefore also give
$h_\Theta\le C(1+\max\{0,h_{\mathrm F}\})$.
\end{remark}

\subsection{A field of definition for homomorphisms}

We now choose the fiber used in the arithmetic proof.  Fix
$s\in\mathbf P^1(\Q)\setminus\mathcal C$, put $M=M_S(s)$ and
$T=S\cup\{q:q\mid M\}$, and choose $y$ above $s$.  Let
$A=\mathcal A_y$ and choose the splitting field $K_0$ from
Lemma~\ref{hei:fiber-field}.  Since $\pi$ is defined over $\Q$ and $s$
is rational, $\pi({}^g y)=g(s)=s$ for every $g\in G_\Q$.
The common-base property of the family therefore makes the conjugates
of $A$ geometrically isogenous.
The descent argument needs all their homomorphisms over one Galois
field.  We can define them without losing the bounded degree or
introducing ramification at further primes.

\begin{lemma}\label{hei:hom-field}
There is a Galois extension $L/\Q$ containing $K_0$, of uniformly
bounded degree, over which all geometric homomorphisms between all
Galois conjugates of $A$ are defined.  It is unramified outside $T$.
\end{lemma}

\begin{proof}
There are boundedly many conjugates because they are defined over
$K_0$.  Let $\mathcal H$ be the direct sum of the geometric Hom groups
indexed by all ordered pairs of these conjugates.  This is a free
abelian group of bounded rank.  Galois transport gives a natural
action of $G_\Q$ on $\mathcal H$, permuting the summands according to
its action on the ordered pairs.  The action of $G_{K_0}$ on $\mathcal H$ has
finite image: a finite set of generators of the Hom groups is defined
over a finite extension.  Finite subgroups of $\mathrm{GL}_r(\Z)$
have bounded order for bounded $r$.  Let $L$ be the fixed field of
the kernel of this action inside $G_{K_0}$.  The transport action of
$G_\Q$, together with normality of $G_{K_0}$, makes that kernel normal
in $G_\Q$.
Thus $L/\Q$ is Galois and its degree is uniformly bounded.

At a place outside $T$, the extension $K_0/\Q$ is unramified and all
the conjugates have good reduction.  For a coefficient prime different
from the residue characteristic, inertia acts trivially on their Tate
modules and therefore on the Hom groups embedded faithfully between
those modules.  Thus $L/K_0$ introduces no ramification there.
\end{proof}

All homomorphisms between conjugates, including the prescribed
$B$-action, are now defined over $L$.  Its degree is uniformly bounded,
and it is unramified outside $T$.  Extending $K_0$ to $L$ preserves
good reduction of $A$ outside the fixed set $S$.  We now bound
$h_{\mathrm F}(A)$ in terms of $M$, beginning with its endomorphism
algebra.

\subsection{The endomorphism dichotomy and descent up to isogeny}
\label{hei:endomorphisms}

The endomorphism algebra determines which comparison will give a
height bound for $A$.  Put
$D=\operatorname{End}^0_{\overline\Q}(A)$, and let $C_B$ be the
centralizer of the prescribed left $B$-action in $D$.  On rational
homology the centralizer of left $B$ is right $B$.  A Hodge
endomorphism in that centralizer must commute, over $\mathbf R$, with
the complex structure, whose centralizer in $M_2(\mathbf R)$ is
$\mathbf C$.  Thus $C_B$ is either $\Q$ or an imaginary quadratic
field $k$ contained in $B$.  More explicitly, a nonscalar rational
element commuting with the complex structure generates such a
quadratic field, and its real centralizer has dimension two, leaving
no further possibilities.

Both cases use the following quantitative isogeny theorem to transfer
a known height bound.  If $C_1,C_2$ are $g$-dimensional abelian
varieties over a number field $k_0$ of degree $d$, and are isogenous
over an extension $K/k_0$, then there are $K$-isogenies in both
directions with degree at most
\begin{equation}\label{hei:gr}
 \kappa(C_1)=
 \left((14g)^{64g^2}d
       \max\{h_{\mathrm F}(C_1),\log d,1\}^{2}\right)^{1024g^3}.
\end{equation}
This is \cite[Theorem 1.4, with the definition on p.~2058]{GaudronRemond2014}.
It uses the original stable Faltings height and requires neither
semistability over $k_0$ nor a principal polarization.  In our
applications we take $K=k_0$.  The stable height is additive on
products and satisfies
\begin{equation}\label{hei:isogeny-height}
 h_{\mathrm F}(C_2)\le h_{\mathrm F}(C_1)
                         +\tfrac12\log\deg\phi
 \quad\text{for an isogeny }\phi:C_1\longrightarrow C_2
\end{equation}
\cite{Faltings1983}.  We will always apply \eqref{hei:gr} to a source
whose height has already been bounded.

If $C_B=k$, the central-simple-algebra double-centralizer calculation
gives
\[
                         D=B\otimes_\Q k\simeq M_2(k).
\]
The last isomorphism holds because an embedded quadratic field is a
maximal subfield of $B$ and splits it.  Matrix idempotents, all defined
over $L$ by Lemma~\ref{hei:hom-field}, give
$A\sim_L E^2$ for a CM elliptic curve $E/L$: take the image of an
integer multiple of a rank-one idempotent, and use the matrix units
to identify the two elliptic factors up to isogeny.  The degree of
$j(E)$ is at most $[L:\Q]$, hence is bounded.  Complex multiplication
theory identifies the degree of a singular modulus with the class
number of its imaginary quadratic order.  Class-number growth for
imaginary quadratic fields and the formula for class numbers of
orders imply that only finitely many such $j$-invariants have bounded
degree \cite{Cox2013,Siegel1935}.  Their stable elliptic heights, and
therefore $h_{\mathrm F}(E^2)$, are bounded.  Apply \eqref{hei:gr} to
the source $E^2$, with dimension two, bounded height, and bounded
$[L:\Q]$, and then apply \eqref{hei:isogeny-height}.  This bounds
$h_{\mathrm F}(A)$ uniformly.  Lemma~\ref{hei:base-height} already
proves Theorem~\ref{hei:height} for these fibers, since $M\ge1$.
The finiteness invoked here need not be effective.

We henceforth treat the other case, $C_B=\Q$.  The same
double-centralizer calculation gives $D=B$.  One can check this
calculation after an extension splitting $B$: an algebra containing
the full $2\times2$ matrix units is a matrix algebra over their
centralizer.  Because $B$ is division, $A$ is absolutely simple; a
proper isogeny factor would give a nontrivial idempotent in $D$.

Our goal in this case is to find a primitive weight-two newform $f$
of level $N$, a Galois extension $L'/\Q$ containing $L$, and an
abelian variety $D_0/L'$ such that, for fixed constants $C,a>0$,
\[
 [L':\Q]\le CM,\qquad N\le CM^a,\qquad
                    J_1(N)\sim_{L'}A\times D_0.
\]
Once the height of this modular Jacobian has been bounded in terms
of $N$, the factorization will let us transfer the bound to $A$ over
a field of controlled degree.  We first use the conjugacy isogenies
to construct a two-dimensional representation.  Their discrepancy
under composition will be reduced to a sign and then split with
control on both the field degree and inertia.

Compatible-isogeny cocycles and their scalar splittings are part of
the theory of $\Q$-curves and building blocks
\cite[Section~6]{Ribet1992}\cite[Chapters~4--5]{Pyle1995};
compare \cite[Section~4]{Quer2001Fields}.
The scalar cochain constructed below takes values among roots of unity
of order $O(M)$, and its inertia restrictions will control the modular level.

Put $\Gamma=\operatorname{Gal}(L/\Q)$.  For each $g\in\Gamma$, choose
a quasi-isogeny $\mu_g:{}^gA\to A$, with $\mu_1=1$, satisfying
\begin{equation}\label{hei:blinear}
             \mu_g\,{}^g\alpha=\alpha\mu_g\qquad(\alpha\in B).
\end{equation}
To obtain it, start with an isogeny between the fibers above the
common rational point.  Transporting the $B$-action along this
isogeny gives an automorphism of $B=D$, which is inner by
Skolem--Noether.  Composing with the corresponding element of
$B^\times$ gives \eqref{hei:blinear}.  All these maps are defined
over $L$ by Lemma~\ref{hei:hom-field}.

The maps $\mu_g\,{}^g\mu_h$ and $\mu_{gh}$ have the same source and
target and both intertwine the $B$-actions.  Their quotient therefore
lies in the centralizer $\Q^\times$:
\[
                 \mu_g\,{}^g\mu_h=c(g,h)\mu_{gh},
                 \qquad c(g,h)\in\Q^\times.
\]
For a quasi-isogeny, use the degree obtained by extending the
geometric degree multiplicatively from isogenies.  A rational scalar
$c$ on a surface has degree $c^4$, so the last identity gives
\[
 |c(g,h)|=
 \left(\frac{\deg\mu_g\,\deg\mu_h}{\deg\mu_{gh}}\right)^{1/4}.
\]
Let $a_g=(\deg\mu_g)^{-1/4}$ be the positive real algebraic fourth
root.  We extend the coefficient scalars in the isogeny category and
put
\[
 \widetilde\mu_g=a_g\mu_g
 \in\operatorname{Hom}_L^0({}^gA,A)\otimes_\Q\overline\Q.
\]
Here Galois transports the Hom factor and acts trivially on the
coefficient factor $\overline\Q$.  The preceding degree identity
now gives
\begin{equation}\label{hei:sign-cocycle}
       \widetilde\mu_g\,{}^g\widetilde\mu_h
              =\epsilon(g,h)\widetilde\mu_{gh},
       \qquad
       \epsilon(g,h)=\frac{c(g,h)}{|c(g,h)|}\in\{\pm1\}.
\end{equation}
Associativity makes $\epsilon$ a normalized cocycle on $\Gamma$;
we inflate it to $G_\Q$.  This is the obstruction in the isogeny
version of Weil descent.  The normalized discrepancy depends only on
the sign of $c(g,h)$, so we will not need bounds for the individual
degrees $\deg\mu_g$.

A scalar cochain with coboundary $\epsilon$ will cancel this sign.
On $G_L$ the inflated cocycle is trivial, so the restriction of such
a cochain will be a character.  We seek one whose restricted character
is conjugacy invariant and has controlled order: its kernel will then
define a Galois field over which the corrected action agrees with the
Tate action of $A$.  Control on inertia will also be needed to bound
the modular level.  The following splitting supplies both properties.

\begin{lemma}[Controlled splitting of the sign obstruction]
\label{hei:sign-splitting}
There is a finite Dirichlet character $\chi$ of order $m\le C M$,
unramified outside $T$, and a continuous cochain
$\beta:G_\Q\to\mu_{2m}$ such that
\begin{equation}\label{hei:beta}
 \beta^2=\chi,\qquad
 \frac{\beta(g)\beta(h)}{\beta(gh)}=\epsilon(g,h).
\end{equation}
The cochain can be chosen trivial on inertia at every prime outside
$T$.  Its restriction to $G_L$ is a conjugacy-invariant character.
The local conductor of $\chi$ at any fixed prime is bounded
independently of $s$, and $\chi$ is tame at odd primes.
\end{lemma}

\begin{proof}
Identify
$H^2(G_\Q,\{\pm1\})=\operatorname{Br}(\Q)[2]$.  At a prime
outside $T$, the cocycle factors through the unramified quotient
$\widehat\Z$, which has cohomological dimension one on torsion
coefficients.  Its local class is therefore zero.  We use local
reciprocity and the Brauer local-global theorem in the form described
in \cite{Serre1979,MilneCFT2020,MilneEtale2013}.

For a finite character $\chi_v$ with values in
$\overline\Q^{\,\times}$ with trivial coefficient action, its
square-root obstruction $\partial\chi_v$ in
$H^2(G_{\Q_v},\{\pm1\})$ vanishes if and only if $\chi_v$ has a
continuous character square root.  By reciprocity this is equivalent
to its being trivial on the order-two subgroup of $\Q_v^\times$,
that is, on $-1$.  To verify the equivalence, use
$\Z_p^\times\simeq\mu_{p-1}\times\Z_p$ for odd $p$, and
$\Z_2^\times\simeq\{\pm1\}\times\Z_2$ at $2$.  A character of the
infinite cyclic valuation factor has a square root by choosing a
square root of its value on a uniformizer; the procyclic free factors
also admit character square roots.  The only obstruction in the
finite cyclic torsion factors is the value on their element of order
two.  At the real place the same assertion follows directly from
$G_{\mathbf R}=\Z/2\Z$.

At each odd finite prime with nonzero local class of $\epsilon$,
choose an odd character modulo that prime, so its value at $-1$ is
$-1$.  At $2$, if necessary, use the nontrivial character modulo
$4$.  Their product is a global Dirichlet character $\chi$.  Its
local square-root obstruction matches that of $\epsilon$ at every
finite place.  The real invariant matches too, because the sum of
local invariants of a Brauer class is zero.  Injectivity in the
Brauer local-global theorem gives $\partial\chi=[\epsilon]$.

Each odd-prime factor has order at most $q-1$, and the factor at $2$
has order two.  Thus the order $m$ of their product satisfies
\[
              m\le 2\prod_{q\in T,\ q\ne2}(q-1)\le C M.
\]
The conductor at an odd prime divides that prime, and the conductor
at $2$ divides $4$.  Factors at other primes are unramified locally;
this proves the stated local conductor assertions.

Choose a continuous set-theoretic square root of $\chi$, valued in
$\mu_{2m}$.  Equality of the obstruction classes allows its
coboundary to be corrected to exactly $\epsilon$ by a sign-valued
continuous cochain.  The resulting $\beta$ satisfies
\eqref{hei:beta}.  On $G_L$ the inflated cocycle is identically one,
so $\beta$ restricts to a character.  Its conjugacy invariance follows
from the cochain identity: if $u\in G_L$ and $g\in G_\Q$, then
$\epsilon(g,u)=\epsilon(gug^{-1},g)=1$, and comparison of
$\beta(gu)$ computed in these two ways gives
$\beta(gug^{-1})=\beta(u)$.

At an odd prime outside $T$, inertia fixes $L$, and $\chi$ is
unramified, so the restriction of $\beta$ to inertia is a sign
character.  Only finitely many such restrictions are nontrivial,
because a continuous function to the finite set $\mu_{2m}$ factors
through a finite quotient of $G_\Q$.  At an odd prime there is a
unique nontrivial quadratic character of inertia: wild inertia is
pro-odd and the tame quotient has a unique quotient of order two.
Let $d_\beta$ be a squarefree integer whose odd prime factors are
precisely these offending primes.  Multiplication of $\beta$ by the
quadratic character of $\Q(\sqrt{d_\beta})$ cancels all the unwanted
inertia characters.  This quadratic character is unramified away
from those primes and possibly $2\in S$.  Its square is one and its
coboundary is one, so \eqref{hei:beta} and the restriction properties
are preserved.  Since $\{\pm1\}\subset\mu_{2m}$, the corrected cochain
still takes values in $\mu_{2m}$, whose elements have orders dividing
$2m\le CM$.
\end{proof}

\subsection{Two-dimensional representations and modularity}
\label{hei:modularity}

We construct representations at coefficient primes $2$ and $3$.
The $2$-adic one will yield a modular form.  The $3$-adic one will
measure the conductor at $2$, where the conductor comparison requires
a coefficient prime different from the residue prime.  This requires
identifying both representations with the same newform.

Fix coefficient embeddings $\overline\Q\hookrightarrow
\overline\Q_\ell$ for $\ell=2,3$.  On the covariant rational Tate
module of $A$ with these extended coefficients put
\begin{equation}\label{hei:representation}
 R_\ell(g)=\beta(g)\widetilde\mu_{\bar g}\circ g,
       \qquad g\in G_\Q,\quad \bar g=g|_L.
\end{equation}
Here $g$ first transports torsion points to \({}^gA\), and the
quasi-isogeny returns them to $A$.  In a composition the discrepancy
of the maps is $\epsilon(g,h)$, and the discrepancy of the scalar
cochain is the same sign.  Their product is one.  Thus $R_\ell$ is
a genuine continuous representation.  All algebraic coefficients
and the splitting of $B$ can be taken in one finite extension of
$\Q_\ell$ for each $\ell$.

Equation \eqref{hei:blinear} shows that $R_\ell$ commutes with $B$.
After coefficient extension, $B\otimes\overline\Q_\ell\simeq
M_2(\overline\Q_\ell)$; the matrix units then give
\begin{equation}\label{hei:double-copy}
                         R_\ell\simeq r_\ell\oplus r_\ell
\end{equation}
for a two-dimensional representation $r_\ell$.  On $G_L$ the
quasi-isogeny in \eqref{hei:representation} is
$\widetilde\mu_1=1$, so only the character $\beta|_{G_L}$ separates
$R_\ell$ from the Tate action of $A$.  Define
\[
                   G_{L'}=\ker(\beta|_{G_L}).
\]
Conjugacy invariance of this character makes $L'/\Q$ Galois, and its
image has order at most $2m$.  Consequently
\begin{equation}\label{hei:lprime-degree}
                     [L':\Q]\le 2m[L:\Q]\le C M.
\end{equation}
On $G_{L'}$, the representation $R_\ell$ is exactly the Tate-module
representation of $A$, with coefficients extended.  This field will
allow us to recover an abelian factor over $L'$.  Its degree may grow
with $M$; the local conductor argument will use $L$, whose degree is
uniformly bounded.

\begin{samepage}
The modularity input is \cite[Theorem 1.1]{OddRegularFontaineMazurAtTwo}
in its exact form: every continuous irreducible odd two-dimensional
$2$-adic representation of $G_\Q$, unramified outside finitely many
primes and de Rham at $2$ with distinct Hodge--Tate weights, is a Tate
twist of the representation associated with a classical cuspidal
eigenform.  There is no residual-image hypothesis in that statement.
We now verify its hypotheses for $r_2$.  The same arguments also give
the properties of $r_3$ used below.  The commutant and
Hodge-decomposition arguments have antecedents in
\cite[Section~3]{Ribet1992}; we apply them to the corrected action
\eqref{hei:representation}.
\end{samepage}

By Faltings' semisimplicity and endomorphism theorems,
\[
 \operatorname{End}_{G_{L'}}(V_\ell(A))
       =\operatorname{End}_{L'}(A)\otimes\Q_\ell
       =B\otimes\Q_\ell
\]
\cite{Faltings1983}.  Therefore $r_\ell|_{G_{L'}}$ is semisimple
with scalar commutant after algebraic coefficient extension, and is
absolutely irreducible.  So is $r_\ell$.

The de Rham comparison theorem for smooth proper varieties, together
with its Hodge--Tate decomposition, applies to $A$ over every
completion of $L'$ at its coefficient prime
\cite[Corollary~1.8]{Scholze2013DeRham}; see also the corrected
definitions in \cite{Scholze2016DeRhamCorrigendum}.
Since the covariant Tate module is dual to first \'etale cohomology,
$V_\ell(A)$ is de Rham with cyclotomic exponents $0,1$, each twice.
The de Rham property descends through finite local extensions
\cite[Section~3.9(ii)]{Fontaine1994}.
Thus $R_\ell$ is de Rham over $\Q_\ell$, and its decomposition
\eqref{hei:double-copy} gives exponents $0,1$, each once, for $r_\ell$.
Here the cyclotomic character has exponent $1$ and Hodge--Tate
weight $-1$.

For complex conjugation $j$, the involution $r_\ell(j)$ cannot be
scalar.  If it were, \eqref{hei:double-copy} would make $R_\ell(j)$
scalar too.  Under comparison with Betti homology,
$\mu_{\bar j}\circ j$ is conjugation followed by a holomorphic
quasi-isogeny.  It exchanges the two nonzero Hodge subspaces, whereas
a scalar preserves both.  Multiplying by the nonzero algebraic
coefficient in \eqref{hei:representation} does not change this
fact.  Hence the two eigenvalues of $r_\ell(j)$ are $1$ and $-1$,
and $r_\ell$ is odd.

At a prime outside $T\cup\{\ell\}$, inertia fixes $L$, $A$ has
good reduction, and the corrected $\beta$ is trivial.  Thus $r_\ell$
is unramified there.  In particular it is ramified at only finitely
many primes.  All hypotheses of the stated modularity theorem now
hold for $r_2$.  Applying it realizes $r_2$ as a Tate twist of the
representation associated with a classical cuspidal eigenform.

We use the arithmetic-Frobenius normalization in which a weight-$k$
form $f$ has good traces $a_q(f)$ and determinant
$\psi_f\chi_{\mathrm{cyc}}^{k-1}$.  Its cyclotomic exponents are
$0,k-1$; the companion fixes this convention in
\cite[Section~2.1]{OddRegularFontaineMazurAtTwo}.  The two exponents
$0,1$ of $r_2$ therefore force weight two and Tate shift zero.  Thus
$r_2$ is realized by the arithmetic $2$-adic representation of a
normalized primitive weight-two newform.

For comparison with a dual presentation, if an intermediate realization is written as
$V_{f,2}^{\vee}\otimes\chi_{\mathrm{cyc}}^n$, then
\[
 V_{f,2}^{\vee}\otimes\chi_{\mathrm{cyc}}^n
 \simeq V_{g,2}\otimes\chi_{\mathrm{cyc}}^{n+1-k},
 \qquad g=(f\otimes\psi_f^{-1})_{\mathrm{new}}.
\]
Here $g$ is the primitive newform of the indicated finite-character
twist.  The target exponents $0,1$ give $k=2$, $n=1$ in this
presentation, and hence arithmetic Tate exponent zero.  Changing
from arithmetic to geometric Frobenius inverts the evaluated
eigenvalues; it does not dualize the representation.  After this
conversion when necessary, let $f$ denote the normalized primitive
newform whose arithmetic $2$-adic representation realizes $r_2$,
and let $N$ be its level.

To use $r_3$ for the dyadic conductor of this level, we must identify
it with the same algebraic conjugate of $f$.  A single algebraic
trace, obtained from the good special fiber, will compare the two
coefficient primes.

\begin{lemma}\label{hei:compatibility}
The representation $r_3$ is the $3$-adic representation of the same
algebraic conjugate of $f$ that realizes $r_2$.
\end{lemma}

\begin{proof}
Fix a rational prime $q\notin T\cup\{2,3\}$ with $q\nmid N$,
a place of $\overline\Q$ above $q$, and an arithmetic Frobenius
$g$ in that decomposition group.  These primes form a cofinite set,
as needed for the Chebotarev argument below.  The automorphism induced
by $g$ on the residue field $\overline{\mathbf F}_q$ is $x\mapsto x^q$.
Therefore the reduction of
\({}^gA\) is the Frobenius twist $\widetilde A^{(q)}$ of the same
geometric special fiber $\widetilde A$.  Specialization of
prime-to-$q$ torsion identifies the action of $g$ with the relative
Frobenius map $\widetilde A\to\widetilde A^{(q)}$.  A multiple of
$\mu_{\bar g}$ extends over good reduction, and its reduction, after
dividing by that integer, maps $\widetilde A^{(q)}$ back to
$\widetilde A$.  The composite is thus an element
\[
       u_q=\operatorname{red}(\mu_{\bar g})\circ
                    \operatorname{Frob}^{\mathrm{rel}}_q
                    \in\operatorname{End}^0(\widetilde A).
\]
Here $\operatorname{red}(\mu_{\bar g})$ is reduction of the
unnormalized quasi-isogeny; the scalar $a_{\bar g}$ from degree
normalization is inserted below.

The characteristic polynomial of an endomorphism of an abelian
variety on rational Tate modules has rational coefficients and is
independent of the coefficient prime.  This applies to a rational
endomorphism after clearing denominators
\cite[I, Theorem 10.9 and Proposition 10.20]{MilneAV2008}.
Consequently
\[
       \alpha_q=\frac12\beta(g)a_{\bar g}\operatorname{Tr}(u_q)
                         \in\overline\Q
\]
is a single algebraic number whose images at $2$ and $3$ are the
traces of $r_2(g)$ and $r_3(g)$, respectively.  The factor $1/2$
comes from \eqref{hei:double-copy}.  This construction is independent
of any comparison of coefficient-field degrees.

The identification of $r_2$ with the representation of $f$ selects
an embedding of its Hecke field in $\overline\Q$.  Injectivity of
the fixed embedding $\overline\Q\hookrightarrow\overline\Q_2$
then gives $\alpha_q=a_q(f)$ as algebraic numbers at every good
prime.  Under the chosen embedding at $3$, the same equalities
hold.  Chebotarev density and the character criterion for semisimple
representations \cite[IV, Theorem~3.3 and Lemma~3.9]{MilneAV2008}
identify $r_3$ with that member of the compatible system of $f$
\cite{Deligne1971Modular}.  Semisimplicity on our side follows already
from irreducibility.
\end{proof}

\subsection{A uniform bound for the modular level}

\begin{lemma}\label{hei:level-bound}
There are fixed constants $C,a>0$ such that
\begin{equation}\label{hei:level}
                              N\le C M^a.
\end{equation}
\end{lemma}

\begin{proof}
For a primitive weight-two form, the local newform conductor is
the exponent of $q$ in its primitive level.  Local--global
compatibility, with conductor preservation under the local
correspondence, identifies this with the Artin-conductor exponent
of its $\ell$-adic representation whenever $q\ne\ell$
\cite[Theorem~A in Section~0.7, with the normalization in Section~0.5]{Carayol1986}.
The conductor includes the Weil--Deligne monodromy contribution.
We use $r_2$ at all primes other than $2$, and $r_3$ at $2$,
justified by Lemma~\ref{hei:compatibility}.  The exponent is zero at
$q\notin T$ by the unramified assertion above.

First let $q\in T$ be larger than a fixed constant that includes all
primes in $S$, the primes $2,3$, and the uniform upper bound for
$[L:\Q]$.  Wild inertia, a pro-$q$ group, acts trivially on $L$.
The surface has good reduction over $L$ at this prime, so wild
inertia also acts trivially on its prime-to-$q$ Tate modules.  On
this subgroup $\beta$ is a character, and its square is the tame
character $\chi$.  Hence its wild restriction is sign-valued and
must be trivial, as $q$ is odd.  Formula \eqref{hei:representation}
makes $R_\ell$ tame at $q$, and hence $r_\ell$ is tame there as well.
Its dimension is two, so its
conductor exponent is at most two.

It remains to bound the exponents at the fixed finite set of smaller
primes.  Fix one of them, say $q$, a place $v$ of $L$ above it, and
put $F=L_v$.  Let $\ell\ne q$ be the coefficient prime chosen above.
Only finitely many embedded extensions $F/\Q_q$ can occur, because
their degrees are bounded.  We seek a bound $u_q$, independent of
$s$, such that for every $u>u_q$ the upper ramification group
$G_{\Q_q}^u$ lies in $G_F$, is killed by the character
$\beta|_{G_F}$, and acts trivially on $V_\ell(A)$.  On such a group
$g|_L=1$, so the quasi-isogeny in $R_\ell(g)$ is also the identity.
We obtain this bound by controlling the character and the Tate action
over the finitely many possible fields $F$.

First consider the character.  By local reciprocity, the square of
$\beta|_{G_F}$ corresponds to
$\chi\circ\operatorname{Norm}_{F/\Q_q}$.  Although the global order
of $\chi$ may grow with $M$, its conductor on $\Q_q$-units is
uniformly bounded: of the constructed Dirichlet factors only the
factor at $q$ ramifies there.  Choose an integer $b$ large enough,
uniformly over the finitely many possibilities for $F$, that this
square character is trivial on $U_F^b=1+\mathfrak m_F^b$ and that
the local logarithm identifies $U_F^b$ with $\mathfrak m_F^b$.
Then $\beta(U_F^b)\subset\{\pm1\}$.  For odd $q$, squaring maps
$U_F^b$ onto itself, and therefore $\beta$ kills $U_F^b$.  For
$q=2$, writing $e_F=v_F(2)$ gives
\[
                         (U_F^b)^2=U_F^{b+e_F}
\]
in these logarithmic coordinates.  Thus $\beta$ kills a uniformly
deeper unit group also at $2$.  Reciprocity gives a uniform upper
ramification cut for this character.  The quadratic correction in
Lemma~\ref{hei:sign-splitting} does not affect the argument, because
it leaves $\beta^2=\chi$ unchanged.  An arbitrarily large unramified
character order has no effect on a unit-filtration bound.

Next consider the Tate action.  Choose a fixed prime $p_0\ge3$,
different from $q$, and put
\[
                           F_1=F(A[p_0]).
\]
Its degree satisfies $[F_1:F]\le|\mathrm{GL}_4(\F_{p_0})|$.
The semistable reduction criterion says that inertia is unipotent
over $F_1$ \cite[Theorem 3.5]{SilverbergZarhin1995}.  Its wild
pro-$q$ subgroup has trivial action on $V_\ell(A)$, since a compact
unipotent subgroup of $\mathrm{GL}_4(\Q_\ell)$ has no nontrivial
pro-$q$ subgroup when $\ell\ne q$.  Since both $[F:\Q_q]$ and
$[F_1:F]$ are bounded, the fields $F_1$ and their normal closures
over $\Q_q$ also range over finite sets.

Choose an upper ramification cut over $\Q_q$ exceeding the breaks
of all these finite field extensions and whose images under the
finitely many relevant Herbrand inverse functions exceed the cuts
for $\beta|_{G_F}$.  The subgroup and quotient rules for upper
ramification groups then place every sufficiently high
$G_{\Q_q}^u$ inside $G_F$, in the kernel of $\beta|_{G_F}$, and in
the wild Tate-module kernel over $F_1$ \cite{Serre1979}.  On this
group $\widetilde\mu_{\bar g}=\widetilde\mu_1=1$, so
\eqref{hei:representation} makes $R_\ell$ trivial there.  Its
decomposition \eqref{hei:double-copy} then makes $r_\ell$ trivial
there as well.  We have therefore obtained the required uniform
upper cut $u_q$ for $r_\ell$, independent of $s$.  Its Swan conductor
is at most $2u_q$.  The remaining contribution from inertia
invariants, including Weil--Deligne monodromy, is at most two because
$r_\ell$ has dimension two.  This bounds its full conductor exponent
at the fixed prime.

The finitely many fixed primes contribute a constant to $N$.  Every
remaining ramified prime divides $M$ and has exponent at most two,
so their product contributes at most $M^2$.  This proves
\eqref{hei:level}, and the proof allows $a=2$ after altering the fixed
constant.  The uniform local bounds use the bounded degree of $L$;
\eqref{hei:lprime-degree} supplies the degree of the field $L'$ for
the global isogeny estimate.
\end{proof}

\subsection{The modular factor and the final height bound}
\label{hei:final-height}

For elliptic curves over \(\mathbb Q\), Murty and Pasten earlier
obtained an effective \(O(N\log N)\) bound for Faltings height in
terms of the conductor, with applications to \(S\)-unit equations
\cite[Theorems~1.1 and~7.1]{MurtyPasten2013}.  The abelian-variety
comparison used in this final step follows the modular height
strategy of von K\"anel \cite[Sections~1.2 and~5]{vonKaenel2021};
see also \cite{Alpoge2021}.
The preceding descent and ramification arguments provide the
point-dependent fields and the level bound needed to apply that strategy
here.

Let $J=J_1(N)$.  The weight-two Eichler--Shimura realization places
$r_2$ in $V_2(J)\otimes\overline\Q_2$
\cite[Theorem~2.1 and the following weight-two realization]{Ribet1977Nebentypus}.
We will turn this shared constituent into a nonzero homomorphism
$A\to J$ over $L'$.  The next argument descends a Hom from the
extended coefficient field to $\Q_2$, giving the abelian factor over
$L'$ itself.

\begin{lemma}\label{hei:modular-factor}
There is an abelian variety $D_0/L'$ such that
\[
                             J\sim_{L'}A\times D_0.
\]
\end{lemma}

\begin{proof}
Put $V=V_2(A)$, $W=V_2(J)$, and $H=G_{L'}$.  After extending
coefficients, $V$ is two copies of $r_2|_H$, while $W$ contains at
least one copy.  Hence
$\operatorname{Hom}_H(V\otimes\overline\Q_2,
W\otimes\overline\Q_2)\ne0$.  Invariants commute with coefficient
extension in this situation.  Indeed, in the finite-dimensional
space $\operatorname{Hom}_{\Q_2}(V,W)$ they are the intersection of
the kernels of the linear equations expressing commutation with
each $h\in H$.  Finitely many of these equations cut out that same
intersection, by stabilization of dimensions.  Kernels of this
finite system commute with scalar extension.  Thus
\[
 \operatorname{Hom}_H(V,W)\otimes\overline\Q_2
 =\operatorname{Hom}_H(V\otimes\overline\Q_2,
                     W\otimes\overline\Q_2)\ne0.
\]
Faltings' homomorphism theorem identifies
$\operatorname{Hom}_H(V,W)$ with
$\operatorname{Hom}_{L'}(A,J)\otimes\Q_2$, so an actual nonzero
$L'$-homomorphism $A\to J$ exists.  Absolute simplicity makes its
kernel finite.  Poincare complete reducibility over $L'$ then gives
the claimed factor \cite{Faltings1983,Mumford1970}.  A single shared
constituent was sufficient to produce the Hom; the resulting rational
Tate-module injection automatically supplies its full required
multiplicity.  In particular $\dim J\ge2$.
\end{proof}

The controlled field, level, and factorization sought in the non-CM
case are now established.  We turn to the height of the modular
Jacobian.  Write $g_0=\dim J$.  The modular group index bounds the
degree of the map $j:X_1(N)\to\mathbf P^1$ by $C N^2$.  It is
branched only over $0,1728,\infty$, so after rescaling the target it
is a Belyi map.  Riemann--Hurwitz, or the usual genus formula, also
gives $g_0\le C N^2$.  Javanpeykar's theorem
\cite[Theorem 1.1.1 and Section 2.3]{Javanpeykar2014} for a smooth
projective connected curve $X$ of genus $g\ge1$ bounds its Jacobian's
stable height in his convention by
\[
                 13\cdot10^6 g\,\deg_B(X)^5.
\]
Javanpeykar's height here equals the original stable Faltings
height of the Jacobian used in \eqref{hei:gr}.  The addition of
$(g/2)\log\pi$ in \cite[Lemma 2.4.4]{Javanpeykar2014} converts
instead to the Deligne normalization $h_{\mathrm F}^{\mathrm D}$
of \cite[Section 2.3]{GaudronRemondPeriods2014}.
Thus no normalization correction is needed here; in either
convention the difference is only $O(g)$.
Since Lemma~\ref{hei:modular-factor} ensures $g_0\ge2$,
the genus hypothesis applies.  We obtain
\begin{equation}\label{hei:jacobian-height}
         g_0\le C N^2,
         \qquad \max\{0,h_{\mathrm F}(J)\}\le C N^{12}.
\end{equation}

For the final comparison we choose a principally polarized complement.
An abelian variety over a number field has a polarization defined over
that field and is consequently isogenous there to its dual.  When
$D_0\ne0$, choose such a polarization over $L'$.  Applying Zarhin's
trick to it gives a principal polarization over $L'$ on
\[
                 D_Z=(D_0\times D_0^\vee)^4
\]
\cite[Theorem~1.1(ii) and Remark~1.4]{Zarhin2023Quaternion}; see also
\cite[proof of Corollary~3.2]{Bost1996Bourbaki}.  If $D_0=0$, omit
this factor and take all its height and dimension contributions below
to be zero.  In either case this choice of complement gives
\[
                 C_1=J^8\sim_{L'} C_2=A^8\times D_Z.
\]
Both sides have dimension $g=8g_0$.  Their common field has degree
$d=[L':\Q]\le C M$, and the source has known stable height
$h_{\mathrm F}(C_1)=8h_{\mathrm F}(J)$.  From
\eqref{hei:level} and \eqref{hei:jacobian-height},
\[
           g\le C M^{2a},\qquad
           \max\{0,h_{\mathrm F}(C_1)\}\le C M^{12a}.
\]
Applying \eqref{hei:gr} to this source gives an isogeny
$\phi:C_1\to C_2$ whose logarithmic degree is at most
\begin{align}
 \log\deg\phi
 &\le1024g^3\bigl(64g^2\log(14g)+\log d
       +2\log\max\{h_{\mathrm F}(C_1),\log d,1\}\bigr)
       \label{hei:log-isogeny}\\
 &\le C M^C.\nonumber
\end{align}
The last inequality follows by substituting the preceding bounds:
every power of $g$ is a fixed power of $M$, while the remaining
logarithms are bounded by $C(1+\log M)\le C'M$ for $M\ge1$.
This is the point at which the explicit dependence on dimension in
the isogeny theorem is needed.  The theorem is applied with source
$J^8$, whose height is already bounded.

By additivity and \eqref{hei:isogeny-height},
\[
             8h_{\mathrm F}(A)+h_{\mathrm F}(D_Z)
        \le h_{\mathrm F}(C_1)+\tfrac12\log\deg\phi
        \le C M^C.
\]
Bost's lower bound applies to every abelian variety over a number field.
For the complement $D_Z$ chosen above, the form in
\cite[Corollary~8.4]{GaudronRemondPeriods2014} gives
\[
 h_{\mathrm F}^{\mathrm D}(D_Z)\ge
 -\tfrac12(\dim D_Z)\log(2\pi).
\]
The normalization identity
$h_{\mathrm F}^{\mathrm D}=h_{\mathrm F}
+\tfrac12(\dim D_Z)\log\pi$
therefore yields
\[
 h_{\mathrm F}(D_Z)\ge
 -\tfrac12(\dim D_Z)\log(2\pi^2).
\]
Since $\dim D_Z\le g\le C M^{2a}$, it follows that
\begin{equation}\label{hei:faltings-final}
                              h_{\mathrm F}(A)\le C M^C.
\end{equation}
Neither a degree bound for the original $\mu_g$ nor a bound for the
degree of a field of coefficients has entered this estimate.

\begin{proof}[Completion of the proof of Theorem~\ref{hei:height}]
For the arbitrary point $s$ fixed above, let $A=\mathcal A_y$ be the
chosen fiber.  The CM case gave a uniform
upper bound for $h_{\mathrm F}(A)$; the other case gave
\eqref{hei:faltings-final}.  Applying Lemma~\ref{hei:base-height}
gives $\hgt(s)\le H M_S(s)^c$ after enlarging fixed constants $H,c$.
All curves, levels, theta data, exceptional primes, and comparison
constants were fixed independently of $s$.  This proves the
theorem.  Since $M_S(s)\ge1$, the exponent $c$ may, if desired,
be increased to a positive integer.
\end{proof}

\begin{remark}[Related height bounds]
Let \(U\) be a nonempty open subscheme of
\(\operatorname{Spec}\mathcal O_K\) for a number field \(K\).
For fixed relative dimension and \([K:\Q]\), related work bounds
stable Faltings height polynomially in the radical of the product
of the absolute discriminant of \(K\) and the norms of the prime
ideals outside \(U\).
Von K\"anel proves such a bound for abelian schemes over \(U\) whose geometric
generic fiber is simple, non-CM, and \(\Q\)-virtual of
\(\mathrm{GL}_2\)-type, with conjugate isogenies compatible with all
geometric endomorphisms
\cite[Proposition~9.9]{vonKaenel2013Preprint}.
Von K\"anel and Kret treat the product-\(\mathrm{GL}_2\)-type class
with \(G_\Q\)-isogenies and also abelian schemes with geometric CM
\cite[Theorem~7.1]{vonKaenelKret2023}.
\end{remark}
